% =====================================================================================
% EJEMPLOS DE LATEX (para copiar y pegar en tus secciones). Este archivo NO forma parte de la tesina.
% Para ver cómo se ve, agrégalo temporalmente en main.tex:   % =====================================================================================
% EJEMPLOS DE LATEX (para copiar y pegar en tus secciones). Este archivo NO forma parte de la tesina.
% Para ver cómo se ve, agrégalo temporalmente en main.tex:   % =====================================================================================
% EJEMPLOS DE LATEX (para copiar y pegar en tus secciones). Este archivo NO forma parte de la tesina.
% Para ver cómo se ve, agrégalo temporalmente en main.tex:   % =====================================================================================
% EJEMPLOS DE LATEX (para copiar y pegar en tus secciones). Este archivo NO forma parte de la tesina.
% Para ver cómo se ve, agrégalo temporalmente en main.tex:   \input{ejemplos/ejemplos_latex}
% Reglas que verifica el sistema de revisión:
%   * Toda tabla y figura lleva \caption y \label, y se cita en el texto con \ref (nunca «Figura 3» a mano).
%   * Toda cita \cite{clave} existe en referencias/referencias.bib.
% =====================================================================================

\section{Ejemplos de LaTeX} \label{sec:ejemplos}

\subsection{Tabla informativa y tabla comparativa}

Una \textbf{tabla informativa} resume \emph{qué es} cada trabajo; sirve para dar un panorama (p. ej. en los
antecedentes). Una \textbf{tabla comparativa} contrasta los trabajos con criterios que importan para tu proyecto e
incluye tu propuesta; sirve para justificarla. Ninguna sustituye la explicación en texto.

\begin{table}[H]
  \centering
  \caption{Ejemplo de tabla informativa de trabajos previos.}
  \label{tab:ejemplo_informativa}
  \begin{tabular}{@{}p{2.6cm}p{3cm}p{3.6cm}p{5cm}@{}}
    \toprule
    \textbf{Autor (año)} & \textbf{Tema} & \textbf{Metodología} & \textbf{Contribución clave} \\
    \midrule
    Pérez (2018)        & Detección de grietas en concreto. & Visión artificial. & Algoritmo con 90\,\% de precisión. \\
    Gómez y López (2019) & Monitoreo de puentes.            & Sensores de vibración. & Sistema de alerta temprana de fallas. \\
    Martínez (2020)      & Rutas logísticas.                & Algoritmos genéticos.  & Redujo 15\,\% los tiempos de entrega. \\
    \bottomrule
  \end{tabular}
\end{table}

\begin{table}[H]
  \centering
  \caption{Ejemplo de tabla comparativa (incluye la propuesta en la última fila).}
  \label{tab:ejemplo_comparativa}
  \begin{tabular}{@{}p{3.4cm}cccc@{}}
    \toprule
    \textbf{Algoritmo / autor} & \textbf{Precisión (\%)} & \textbf{Velocidad (ms)} & \textbf{Escalabilidad} & \textbf{Hardware} \\
    \midrule
    Pérez (2018)         & 90 & 50  & Baja  & PC estándar \\
    Gómez y López (2019) & 95 & 200 & Media & Servidor \\
    \textbf{Propuesta}   & \textbf{98} & \textbf{35} & \textbf{Alta} & \textbf{PC estándar} \\
    \bottomrule
  \end{tabular}
\end{table}

Como se observa en la Tabla~\ref{tab:ejemplo_comparativa}, la propuesta mejora la precisión sin requerir hardware especial.
Observa que la tabla se cita con \verb|Tabla~\ref{...}| y que el texto la interpreta.

\subsection{Figuras}

\begin{figure}[H]
  \centering
  \includegraphics[width=0.7\textwidth]{img/diagrama_proceso}
  \caption{Diagrama de un proceso de control automático.}
  \label{fig:ejemplo_diagrama}
\end{figure}

La Figura~\ref{fig:ejemplo_diagrama} ilustra el uso de \verb|\includegraphics|: indica la ruta de la imagen, un ancho
relativo al texto, un pie de figura y una etiqueta. Prefiere imágenes vectoriales (PDF) o PNG de buena resolución.

\begin{figure}[H]
  \centering
  \begin{subfigure}[b]{0.45\textwidth}
    \includegraphics[width=\textwidth]{img/grafico_datos1}
    \caption{Datos de temperatura.}
    \label{fig:ejemplo_sub_temp}
  \end{subfigure}
  \hfill
  \begin{subfigure}[b]{0.45\textwidth}
    \includegraphics[width=\textwidth]{img/grafico_datos2}
    \caption{Datos de presión.}
    \label{fig:ejemplo_sub_pres}
  \end{subfigure}
  \caption{Comparación de datos operativos (temperatura y presión).}
  \label{fig:ejemplo_comparacion}
\end{figure}

La Figura~\ref{fig:ejemplo_comparacion} combina dos imágenes con el paquete \verb|subcaption|.

\subsection{Ecuaciones, algoritmos y código}

Una ecuación numerada que se cita en el texto:
\begin{equation}
  \label{eq:ejemplo_error}
  E = \frac{1}{n}\sum_{i=1}^{n}\left(y_i - \hat{y}_i\right)^2 ,
\end{equation}
donde $y_i$ es el valor real y $\hat{y}_i$ el valor estimado (la Ecuación~\ref{eq:ejemplo_error} es el error cuadrático
medio). Define siempre los símbolos que uses.

\begin{algorithm}[H]
  \caption{Búsqueda lineal.}
  \label{alg:ejemplo_busqueda}
  \begin{algorithmic}[1]
    \Require Lista $A$ de $n$ elementos y valor buscado $x$
    \Ensure Posición de $x$ en $A$ o $-1$
    \For{$i \gets 1$ \textbf{to} $n$}
      \If{$A[i] = x$} \State \Return $i$ \EndIf
    \EndFor
    \State \Return $-1$
  \end{algorithmic}
\end{algorithm}

\subsection{Citas}

Con biblatex se cita con \verb|\cite{clave}| (estilo numérico IEEE) o \verb|\textcite{clave}| y \verb|\parencite{clave}|
(estilo APA). La clave debe existir en \texttt{referencias/referencias.bib}.

% Reglas que verifica el sistema de revisión:
%   * Toda tabla y figura lleva \caption y \label, y se cita en el texto con \ref (nunca «Figura 3» a mano).
%   * Toda cita \cite{clave} existe en referencias/referencias.bib.
% =====================================================================================

\section{Ejemplos de LaTeX} \label{sec:ejemplos}

\subsection{Tabla informativa y tabla comparativa}

Una \textbf{tabla informativa} resume \emph{qué es} cada trabajo; sirve para dar un panorama (p. ej. en los
antecedentes). Una \textbf{tabla comparativa} contrasta los trabajos con criterios que importan para tu proyecto e
incluye tu propuesta; sirve para justificarla. Ninguna sustituye la explicación en texto.

\begin{table}[H]
  \centering
  \caption{Ejemplo de tabla informativa de trabajos previos.}
  \label{tab:ejemplo_informativa}
  \begin{tabular}{@{}p{2.6cm}p{3cm}p{3.6cm}p{5cm}@{}}
    \toprule
    \textbf{Autor (año)} & \textbf{Tema} & \textbf{Metodología} & \textbf{Contribución clave} \\
    \midrule
    Pérez (2018)        & Detección de grietas en concreto. & Visión artificial. & Algoritmo con 90\,\% de precisión. \\
    Gómez y López (2019) & Monitoreo de puentes.            & Sensores de vibración. & Sistema de alerta temprana de fallas. \\
    Martínez (2020)      & Rutas logísticas.                & Algoritmos genéticos.  & Redujo 15\,\% los tiempos de entrega. \\
    \bottomrule
  \end{tabular}
\end{table}

\begin{table}[H]
  \centering
  \caption{Ejemplo de tabla comparativa (incluye la propuesta en la última fila).}
  \label{tab:ejemplo_comparativa}
  \begin{tabular}{@{}p{3.4cm}cccc@{}}
    \toprule
    \textbf{Algoritmo / autor} & \textbf{Precisión (\%)} & \textbf{Velocidad (ms)} & \textbf{Escalabilidad} & \textbf{Hardware} \\
    \midrule
    Pérez (2018)         & 90 & 50  & Baja  & PC estándar \\
    Gómez y López (2019) & 95 & 200 & Media & Servidor \\
    \textbf{Propuesta}   & \textbf{98} & \textbf{35} & \textbf{Alta} & \textbf{PC estándar} \\
    \bottomrule
  \end{tabular}
\end{table}

Como se observa en la Tabla~\ref{tab:ejemplo_comparativa}, la propuesta mejora la precisión sin requerir hardware especial.
Observa que la tabla se cita con \verb|Tabla~\ref{...}| y que el texto la interpreta.

\subsection{Figuras}

\begin{figure}[H]
  \centering
  \includegraphics[width=0.7\textwidth]{img/diagrama_proceso}
  \caption{Diagrama de un proceso de control automático.}
  \label{fig:ejemplo_diagrama}
\end{figure}

La Figura~\ref{fig:ejemplo_diagrama} ilustra el uso de \verb|\includegraphics|: indica la ruta de la imagen, un ancho
relativo al texto, un pie de figura y una etiqueta. Prefiere imágenes vectoriales (PDF) o PNG de buena resolución.

\begin{figure}[H]
  \centering
  \begin{subfigure}[b]{0.45\textwidth}
    \includegraphics[width=\textwidth]{img/grafico_datos1}
    \caption{Datos de temperatura.}
    \label{fig:ejemplo_sub_temp}
  \end{subfigure}
  \hfill
  \begin{subfigure}[b]{0.45\textwidth}
    \includegraphics[width=\textwidth]{img/grafico_datos2}
    \caption{Datos de presión.}
    \label{fig:ejemplo_sub_pres}
  \end{subfigure}
  \caption{Comparación de datos operativos (temperatura y presión).}
  \label{fig:ejemplo_comparacion}
\end{figure}

La Figura~\ref{fig:ejemplo_comparacion} combina dos imágenes con el paquete \verb|subcaption|.

\subsection{Ecuaciones, algoritmos y código}

Una ecuación numerada que se cita en el texto:
\begin{equation}
  \label{eq:ejemplo_error}
  E = \frac{1}{n}\sum_{i=1}^{n}\left(y_i - \hat{y}_i\right)^2 ,
\end{equation}
donde $y_i$ es el valor real y $\hat{y}_i$ el valor estimado (la Ecuación~\ref{eq:ejemplo_error} es el error cuadrático
medio). Define siempre los símbolos que uses.

\begin{algorithm}[H]
  \caption{Búsqueda lineal.}
  \label{alg:ejemplo_busqueda}
  \begin{algorithmic}[1]
    \Require Lista $A$ de $n$ elementos y valor buscado $x$
    \Ensure Posición de $x$ en $A$ o $-1$
    \For{$i \gets 1$ \textbf{to} $n$}
      \If{$A[i] = x$} \State \Return $i$ \EndIf
    \EndFor
    \State \Return $-1$
  \end{algorithmic}
\end{algorithm}

\subsection{Citas}

Con biblatex se cita con \verb|\cite{clave}| (estilo numérico IEEE) o \verb|\textcite{clave}| y \verb|\parencite{clave}|
(estilo APA). La clave debe existir en \texttt{referencias/referencias.bib}.

% Reglas que verifica el sistema de revisión:
%   * Toda tabla y figura lleva \caption y \label, y se cita en el texto con \ref (nunca «Figura 3» a mano).
%   * Toda cita \cite{clave} existe en referencias/referencias.bib.
% =====================================================================================

\section{Ejemplos de LaTeX} \label{sec:ejemplos}

\subsection{Tabla informativa y tabla comparativa}

Una \textbf{tabla informativa} resume \emph{qué es} cada trabajo; sirve para dar un panorama (p. ej. en los
antecedentes). Una \textbf{tabla comparativa} contrasta los trabajos con criterios que importan para tu proyecto e
incluye tu propuesta; sirve para justificarla. Ninguna sustituye la explicación en texto.

\begin{table}[H]
  \centering
  \caption{Ejemplo de tabla informativa de trabajos previos.}
  \label{tab:ejemplo_informativa}
  \begin{tabular}{@{}p{2.6cm}p{3cm}p{3.6cm}p{5cm}@{}}
    \toprule
    \textbf{Autor (año)} & \textbf{Tema} & \textbf{Metodología} & \textbf{Contribución clave} \\
    \midrule
    Pérez (2018)        & Detección de grietas en concreto. & Visión artificial. & Algoritmo con 90\,\% de precisión. \\
    Gómez y López (2019) & Monitoreo de puentes.            & Sensores de vibración. & Sistema de alerta temprana de fallas. \\
    Martínez (2020)      & Rutas logísticas.                & Algoritmos genéticos.  & Redujo 15\,\% los tiempos de entrega. \\
    \bottomrule
  \end{tabular}
\end{table}

\begin{table}[H]
  \centering
  \caption{Ejemplo de tabla comparativa (incluye la propuesta en la última fila).}
  \label{tab:ejemplo_comparativa}
  \begin{tabular}{@{}p{3.4cm}cccc@{}}
    \toprule
    \textbf{Algoritmo / autor} & \textbf{Precisión (\%)} & \textbf{Velocidad (ms)} & \textbf{Escalabilidad} & \textbf{Hardware} \\
    \midrule
    Pérez (2018)         & 90 & 50  & Baja  & PC estándar \\
    Gómez y López (2019) & 95 & 200 & Media & Servidor \\
    \textbf{Propuesta}   & \textbf{98} & \textbf{35} & \textbf{Alta} & \textbf{PC estándar} \\
    \bottomrule
  \end{tabular}
\end{table}

Como se observa en la Tabla~\ref{tab:ejemplo_comparativa}, la propuesta mejora la precisión sin requerir hardware especial.
Observa que la tabla se cita con \verb|Tabla~\ref{...}| y que el texto la interpreta.

\subsection{Figuras}

\begin{figure}[H]
  \centering
  \includegraphics[width=0.7\textwidth]{img/diagrama_proceso}
  \caption{Diagrama de un proceso de control automático.}
  \label{fig:ejemplo_diagrama}
\end{figure}

La Figura~\ref{fig:ejemplo_diagrama} ilustra el uso de \verb|\includegraphics|: indica la ruta de la imagen, un ancho
relativo al texto, un pie de figura y una etiqueta. Prefiere imágenes vectoriales (PDF) o PNG de buena resolución.

\begin{figure}[H]
  \centering
  \begin{subfigure}[b]{0.45\textwidth}
    \includegraphics[width=\textwidth]{img/grafico_datos1}
    \caption{Datos de temperatura.}
    \label{fig:ejemplo_sub_temp}
  \end{subfigure}
  \hfill
  \begin{subfigure}[b]{0.45\textwidth}
    \includegraphics[width=\textwidth]{img/grafico_datos2}
    \caption{Datos de presión.}
    \label{fig:ejemplo_sub_pres}
  \end{subfigure}
  \caption{Comparación de datos operativos (temperatura y presión).}
  \label{fig:ejemplo_comparacion}
\end{figure}

La Figura~\ref{fig:ejemplo_comparacion} combina dos imágenes con el paquete \verb|subcaption|.

\subsection{Ecuaciones, algoritmos y código}

Una ecuación numerada que se cita en el texto:
\begin{equation}
  \label{eq:ejemplo_error}
  E = \frac{1}{n}\sum_{i=1}^{n}\left(y_i - \hat{y}_i\right)^2 ,
\end{equation}
donde $y_i$ es el valor real y $\hat{y}_i$ el valor estimado (la Ecuación~\ref{eq:ejemplo_error} es el error cuadrático
medio). Define siempre los símbolos que uses.

\begin{algorithm}[H]
  \caption{Búsqueda lineal.}
  \label{alg:ejemplo_busqueda}
  \begin{algorithmic}[1]
    \Require Lista $A$ de $n$ elementos y valor buscado $x$
    \Ensure Posición de $x$ en $A$ o $-1$
    \For{$i \gets 1$ \textbf{to} $n$}
      \If{$A[i] = x$} \State \Return $i$ \EndIf
    \EndFor
    \State \Return $-1$
  \end{algorithmic}
\end{algorithm}

\subsection{Citas}

Con biblatex se cita con \verb|\cite{clave}| (estilo numérico IEEE) o \verb|\textcite{clave}| y \verb|\parencite{clave}|
(estilo APA). La clave debe existir en \texttt{referencias/referencias.bib}.

% Reglas que verifica el sistema de revisión:
%   * Toda tabla y figura lleva \caption y \label, y se cita en el texto con \ref (nunca «Figura 3» a mano).
%   * Toda cita \cite{clave} existe en referencias/referencias.bib.
% =====================================================================================

\section{Ejemplos de LaTeX} \label{sec:ejemplos}

\subsection{Tabla informativa y tabla comparativa}

Una \textbf{tabla informativa} resume \emph{qué es} cada trabajo; sirve para dar un panorama (p. ej. en los
antecedentes). Una \textbf{tabla comparativa} contrasta los trabajos con criterios que importan para tu proyecto e
incluye tu propuesta; sirve para justificarla. Ninguna sustituye la explicación en texto.

\begin{table}[H]
  \centering
  \caption{Ejemplo de tabla informativa de trabajos previos.}
  \label{tab:ejemplo_informativa}
  \begin{tabular}{@{}p{2.6cm}p{3cm}p{3.6cm}p{5cm}@{}}
    \toprule
    \textbf{Autor (año)} & \textbf{Tema} & \textbf{Metodología} & \textbf{Contribución clave} \\
    \midrule
    Pérez (2018)        & Detección de grietas en concreto. & Visión artificial. & Algoritmo con 90\,\% de precisión. \\
    Gómez y López (2019) & Monitoreo de puentes.            & Sensores de vibración. & Sistema de alerta temprana de fallas. \\
    Martínez (2020)      & Rutas logísticas.                & Algoritmos genéticos.  & Redujo 15\,\% los tiempos de entrega. \\
    \bottomrule
  \end{tabular}
\end{table}

\begin{table}[H]
  \centering
  \caption{Ejemplo de tabla comparativa (incluye la propuesta en la última fila).}
  \label{tab:ejemplo_comparativa}
  \begin{tabular}{@{}p{3.4cm}cccc@{}}
    \toprule
    \textbf{Algoritmo / autor} & \textbf{Precisión (\%)} & \textbf{Velocidad (ms)} & \textbf{Escalabilidad} & \textbf{Hardware} \\
    \midrule
    Pérez (2018)         & 90 & 50  & Baja  & PC estándar \\
    Gómez y López (2019) & 95 & 200 & Media & Servidor \\
    \textbf{Propuesta}   & \textbf{98} & \textbf{35} & \textbf{Alta} & \textbf{PC estándar} \\
    \bottomrule
  \end{tabular}
\end{table}

Como se observa en la Tabla~\ref{tab:ejemplo_comparativa}, la propuesta mejora la precisión sin requerir hardware especial.
Observa que la tabla se cita con \verb|Tabla~\ref{...}| y que el texto la interpreta.

\subsection{Figuras}

\begin{figure}[H]
  \centering
  \includegraphics[width=0.7\textwidth]{img/diagrama_proceso}
  \caption{Diagrama de un proceso de control automático.}
  \label{fig:ejemplo_diagrama}
\end{figure}

La Figura~\ref{fig:ejemplo_diagrama} ilustra el uso de \verb|\includegraphics|: indica la ruta de la imagen, un ancho
relativo al texto, un pie de figura y una etiqueta. Prefiere imágenes vectoriales (PDF) o PNG de buena resolución.

\begin{figure}[H]
  \centering
  \begin{subfigure}[b]{0.45\textwidth}
    \includegraphics[width=\textwidth]{img/grafico_datos1}
    \caption{Datos de temperatura.}
    \label{fig:ejemplo_sub_temp}
  \end{subfigure}
  \hfill
  \begin{subfigure}[b]{0.45\textwidth}
    \includegraphics[width=\textwidth]{img/grafico_datos2}
    \caption{Datos de presión.}
    \label{fig:ejemplo_sub_pres}
  \end{subfigure}
  \caption{Comparación de datos operativos (temperatura y presión).}
  \label{fig:ejemplo_comparacion}
\end{figure}

La Figura~\ref{fig:ejemplo_comparacion} combina dos imágenes con el paquete \verb|subcaption|.

\subsection{Ecuaciones, algoritmos y código}

Una ecuación numerada que se cita en el texto:
\begin{equation}
  \label{eq:ejemplo_error}
  E = \frac{1}{n}\sum_{i=1}^{n}\left(y_i - \hat{y}_i\right)^2 ,
\end{equation}
donde $y_i$ es el valor real y $\hat{y}_i$ el valor estimado (la Ecuación~\ref{eq:ejemplo_error} es el error cuadrático
medio). Define siempre los símbolos que uses.

\begin{algorithm}[H]
  \caption{Búsqueda lineal.}
  \label{alg:ejemplo_busqueda}
  \begin{algorithmic}[1]
    \Require Lista $A$ de $n$ elementos y valor buscado $x$
    \Ensure Posición de $x$ en $A$ o $-1$
    \For{$i \gets 1$ \textbf{to} $n$}
      \If{$A[i] = x$} \State \Return $i$ \EndIf
    \EndFor
    \State \Return $-1$
  \end{algorithmic}
\end{algorithm}

\subsection{Citas}

Con biblatex se cita con \verb|\cite{clave}| (estilo numérico IEEE) o \verb|\textcite{clave}| y \verb|\parencite{clave}|
(estilo APA). La clave debe existir en \texttt{referencias/referencias.bib}.
